% Provenance: paper_P1_scaling (demoted from "law" to observations),
% MATH-500/MBPP panels, frontier synthesis sections, pass@k protocol.
\chapter{Scaling and Transfer Observations}
\label{ch:scaling}

\section{Scope discipline}
The programme's cross-library corpus (70+ logged runs, seven libraries, five
model families, 0.6B--671B) supports \emph{observations}, not laws: most cells
are single-seed, backends differ, and the term ``scaling law'' is deliberately
avoided (external review, Chapter~9). The full scaling-law treatment of the
corresponding working paper is demoted to this observational chapter.

\section{What the corpus shows}
Model scale does not uniformly reduce \ZVF{} ($r\approx-0.26$ across logged
runs), and after accounting for task type, model family alone does not explain
performance variance (one-way ANOVA $F(4,10){=}0.949$, $p{=}0.476$). The
practical reading: signal starvation is a property of the (task difficulty
$\times$ group size $\times$ training phase) geometry, not something scale
buys you out of.

\section{Transfer and difficulty extension}
Post-RL adapters trained on GSM8K show \emph{zero forgetting and mild positive
transfer} on MBPP (pass@32, all adapters within noise of or above base across
$G\in\{2,4,16,32\}$). On the matched base-versus-post-RL pass@$k$ panel,
pass@1 deltas scatter from $-0.5$ to $+1.9$ points --- under pass@1-only
reporting the $G{=}2$ arm would read as a regression --- while the pass@32
frontier improves $+1.5$ to $+2.5$ points in all four arms; none of this
reaches significance at a single training seed, which is exactly why pass@$k$
curves with stated noise bands are a reporting-standard item (Chapter~9).
On MATH-500, the harder-task extension, the base-model panel is complete and
the post-RL panels are partial at the time of writing; the partial result is
that \textbf{GSM8K-trained gains do not replicate on hard math} --- the
frontier shift observed on GSM8K-family tasks does not carry, reinforcing the
distribution-sharpening (rather than capability-expansion) interpretation and
the thesis's non-claim~4.

\section{Takeaway}
Scaling observations here serve one purpose: bounding the claims. The
capability headroom analysis (pass@32 $=91\%$ on GSM8K), the non-transfer to
MATH-500, and the single-seed density of the cross-scale corpus jointly
justify why Claims~1--2 are stated at the stack level and nowhere above it.
